% !PRINTED 731
\weberpage{731}

\noindent
\begin{minipage}[t]{0.48\textwidth}
\noindent Teilungskörper 573, 592, 620.\\
--- und Klassenkörper 620.\\
Thetafunktion 68 ff.\\
Thetanullwerte 85.\\
Thetaquotienten 88 ff.\\
Transformation von Differentialen 5.\\
---, \gesperrt{Gausssche} 33.\\
---, \gesperrt{Hermitesche} 17.\\
---, \gesperrt{Jacobische} 32.\\
---, \gesperrt{Landensche} 34.\\
--- 2. Grades 32.\\
--- 3. Grades 35.\\
--- der $T$-Funktionen 94.\\
Transformationsgleichung 221.\\
---, Bildung der 225 ff.\\
--- erster Stufe 246.\\
--- für $\gamma_2$ und $\gamma_3$ 247 ff.\\
Transformationsgrad 94, 234.\\
Transformationszahl 94, 234.

\noindent \textbf{U}neigentliche Äquivalenz 331.\\
--- Diskriminanten 321.\\
--- Polygonklassen 674.\\
Unstetigkeitspunkte 53.
\end{minipage}%
\hfill\vrule width 0.4pt\hfill
\begin{minipage}[t]{0.48\textwidth}
\noindent \textbf{V}ertauschung von Argument und Parameter 157.\\
Verwandte $T$-Funktionen 65.\\
--- $\theta$-Funktionen 70.\\
Verzweigungspolygon 669.\\
Verzweigungspunkt 669.\\
Verzweigungszahl 669.

\noindent \gesperrt{\textbf{W}eierstrass} 13.\\
\gesperrt{Weierstrasssche} Primfaktoren 56.\\
--- $\wp$-Funktion 150.\\
--- $\sigma$-Funktionen 116.\\
--- Transzendenten 2. und 3. Gattung\\
160 ff.\\
Wurzeln der Transformationsgleichung\\
231.\\
Wurzelsystem, vollständiges 533.

\noindent \textbf{Z}ahlgruppen 358, 596.\\
Zahlringe 351.\\
Zerfällung der Klassengleichung 616.\\
Zusammensetzung der Transformationen 95 ff.\\
Zweiteilung der elliptischen Funktionen\\
200.
\end{minipage}

\begin{center}
\rule{3cm}{0.4pt}
\end{center}
